% answers_short.tex —— 2005 年·普通高等学校招生考试·浙江卷·文科 · 简短答案速查
% 由 tools/make_short_answers.py 自 answers.tex 的 \gksolhead/\gkanswer 提取，
% 并与 2005/2005_数学真题_浙江文科_答案要点.md 逐题比对后生成，勿手改（改源 MD 后重跑）。

\begin{gknotice}
\textbf{使用说明}：本节只给最终答案便于快速核对，完整推导见解析版。答案为编写者独立推导，\textbf{非官方标准答案}。
\end{gknotice}

\gkpart{选择题}
\begin{center}\begin{tabular}{ccccccccccc}
\toprule
题号 & 1 & 2 & 3 & 4 & 5 & 6 & 7 & 8 & 9 & 10 \\
\midrule
答案 & B & A & D & D & D & A & D & C & B & A \\
\bottomrule
\end{tabular}\end{center}

\gkpart{填空题}
\begin{itemize}[leftmargin=2.6em]
  \item[\textbf{11.}] $y = \dfrac{2x}{1 - x}$（$x \neq 1$）。
  \item[\textbf{12.}] $45^\circ$。
  \item[\textbf{13.}] $2$。
  \item[\textbf{14.}] $5832$。
\end{itemize}

\gkpart{解答题}
\begin{itemize}[leftmargin=2.6em]
  \item[\textbf{15.}] (1) $1$；(2) $\dfrac{\sqrt{6} + \sqrt{2}}{4}$
  \item[\textbf{16.}] $a = 11$，$b = 5$，$c = -1$，或 $a = 2$，$b = 5$，$c = 8$。
  \item[\textbf{17.}] (1) ① $\dfrac{40}{243}$；② $\dfrac{1}{27}$；(2) $p = \dfrac{13}{30}$
  \item[\textbf{18.}] (1) $O$、$D$ 为中点，故 $OD \parallel PA$，从而 $OD \parallel$ 平面 $PAB$；(2) $\arcsin\dfrac{\sqrt{210}}{30}$
  \item[\textbf{19.}] (1) $\dfrac{x^2}{4} + \dfrac{y^2}{3} = 1$；(2) $\arcsin\dfrac{1}{4}$
  \item[\textbf{20.}] (1) $g(x) = -x^2 + 2x$；(2) $\left[-1, \dfrac{1}{2}\right]$；(3) $\lambda \leqslant 0$
\end{itemize}
